% TA-PFC-E1 · NEXORA · Paca Pilataza L. A., Velez Vinces J. G. (Paralelo B)
% Compilar: pdflatex main -> bibtex main -> pdflatex main -> pdflatex main
% Las referencias estan en referencias.bib (misma carpeta).
\documentclass[11pt,a4paper]{article}

\usepackage[utf8]{inputenc}
\usepackage[T1]{fontenc}
\usepackage[spanish,es-tabla]{babel}
\usepackage{mathptmx}
\usepackage[margin=2.5cm, headheight=15pt]{geometry}
\usepackage{array}
\usepackage{tabularx}
\usepackage{booktabs}
\usepackage{xcolor}
\usepackage{colortbl}
\usepackage{enumitem}
\usepackage{float}
\usepackage{fancyhdr}
\usepackage{caption}
\usepackage{tikz}
\usetikzlibrary{shapes.geometric, arrows.meta, positioning, calc}
\usepackage[hidelinks]{hyperref}

\definecolor{uteqgreen}{RGB}{18, 112, 58}
\definecolor{tableheader}{RGB}{235, 242, 238}

\newcolumntype{L}[1]{>{\raggedright\arraybackslash}p{#1}}
\newcolumntype{C}[1]{>{\centering\arraybackslash}p{#1}}
\newcolumntype{Y}{>{\raggedright\arraybackslash}X}

\captionsetup{font=small, labelfont=bf, labelsep=period, skip=4pt}
\setlength{\parindent}{0pt}
\setlength{\parskip}{5pt}
\setlength{\emergencystretch}{2em}
\renewcommand{\arraystretch}{1.22}

\pagestyle{fancy}
\fancyhf{}
\fancyhead[L]{\small UTEQ $\cdot$ Ingenier\'ia de Requerimientos (ISR-401)}
\fancyhead[R]{\small TA-PFC-E1 $\cdot$ Planificaci\'on y Elicitaci\'on}
\fancyfoot[L]{\small Paca Pilataza L. A., V\'elez Vinces J. G.}
\fancyfoot[R]{\small P\'agina \thepage}
\renewcommand{\headrulewidth}{0.4pt}
\renewcommand{\footrulewidth}{0.4pt}

\begin{document}

% ==================== ENCABEZADO (ESTILO SGRA) ====================
\begin{center}
    {\Large \textbf{Planificaci\'on y Elicitaci\'on de Requisitos para una Plataforma de Promoci\'on y Difusi\'on de Servicios de Profesionales Independientes y Establecimientos Comerciales}} \\[8pt]
    
    {\normalsize Luis Alberto Paca Pilataza\textsuperscript{1}, Jostin Gustavo V\'elez Vinces\textsuperscript{1}} \\[4pt]
    {\small \textsuperscript{1} Carrera de Ingenier\'ia de Software -- Cuarto Nivel Paralelo B, Facultad de Ciencias de la Computaci\'on, \\
    Universidad T\'ecnica Estatal de Quevedo (UTEQ). Quevedo, Ecuador \\
    \texttt{lpacap@uteq.edu.ec}, \texttt{jvelezv10@uteq.edu.ec}}
\end{center}

\vspace{4pt}

\noindent\textbf{Resumen.} El presente documento establece la base de planificaci\'on y elicitaci\'on de requisitos para el desarrollo de una plataforma orientada a la promoci\'on y difusi\'on de servicios ofrecidos por profesionales independientes y establecimientos comerciales en el cant\'on Quevedo. En la actualidad, la b\'usqueda de servicios como electricidad, plomer\'ia, cerrajer\'ia, carpinter\'ia o atenci\'on en locales comerciales se realiza de manera informal mediante recomendaciones de conocidos o publicaciones sueltas en redes sociales, lo que genera desconfianza, falta de informaci\'on previa y poca visibilidad para los proveedores locales. Con base en los lineamientos de la norma ISO/IEC/IEEE 29148:2018 \cite{iso29148}, Wiegers y Beatty \cite{wiegers2013} y Pohl \cite{pohl2025}, se presenta la ficha del proyecto, la identificaci\'on de las partes interesadas, la delimitaci\'on del alcance del sistema, el diagrama de contexto y el plan de elicitaci\'on de requisitos.

\noindent\textbf{Palabras clave:} Promoci\'on de servicios; profesionales independientes; establecimientos comerciales; partes interesadas; alcance del sistema; elicitaci\'on de requisitos.

\vspace{4pt}
\hrule
\vspace{6pt}

% ==================== A. FICHA DEL PROYECTO ====================
\section*{A. Ficha del proyecto}
\addcontentsline{toc}{section}{A. Ficha del proyecto}

En el cant\'on Quevedo, las personas que necesitan contratar un servicio para el hogar o localizar un negocio específico suelen enfrentar dificultades para encontrar opciones confiables y cercanas a su sector. Al mismo tiempo, los trabajadores independientes y propietarios de locales comerciales no cuentan con un espacio digital centralizado donde puedan mostrar sus servicios, fotograf\'ias de trabajos anteriores, ubicaci\'on y datos de contacto. En la Tabla~\ref{tab:ficha} se presenta la ficha general del proyecto que resume el sistema a desarrollar, el entorno cliente, la problem\'atica y el objetivo general planteado \cite{iso29148}.

\begin{table}[H]
\centering
\caption{Ficha general del proyecto}
\label{tab:ficha}
\small
\begin{tabularx}{\textwidth}{|L{3.8cm}|Y|}
\hline
\rowcolor{tableheader}
\textbf{Elemento} & \textbf{Descripci\'on} \\
\hline
\textbf{Nombre del sistema} & \textbf{NEXORA} --- Plataforma para la promoci\'on y difusi\'on de servicios ofrecidos por profesionales independientes y establecimientos comerciales. \\
\hline
\textbf{Organizaci\'on cliente} & \textbf{Balanzas Electr\'onica CAS} (cant\'on Quevedo, provincia de Los R\'ios), empresa dedicada a la fabricaci\'on, venta y mantenimiento de balanzas y sistemas de pesaje. Particip\'o en el levantamiento de requisitos mediante carta de aceptaci\'on firmada el 20 de noviembre de 2026, en la que autoriz\'o el uso de su nombre en los documentos del proyecto. \\
\hline
\textbf{Usuarios destino} & Clientes o usuarios, profesionales independientes de oficios t\'ecnicos y propietarios de establecimientos comerciales de servicios del cant\'on Quevedo, provincia de Los R\'ios. \\
\hline
\textbf{Limitaci\'on declarada} & El propietario de la organizaci\'on cliente (participante Nexora-01) es familiar directo de un integrante del equipo. Para reducir el sesgo, su entrevista sigui\'o el mismo guion que las dem\'as, se grab\'o con consentimiento firmado y sus hallazgos se contrastan con los de participantes sin relaci\'on con el equipo (Nexora-02 y Nexora-03). \\
\hline
\textbf{Problema identificado} & Actualmente, la b\'usqueda y promoci\'on de servicios en Quevedo se realiza de manera dispersa por medio de recomendaciones verbales, grupos de redes sociales o recorridos presenciales. Los clientes tienen dificultades para encontrar proveedores en su propio sector, desconocen horarios o n\'umeros de contacto actualizados y sienten desconfianza al no poder ver fotograf\'ias de trabajos realizados anteriormente. Por su parte, los profesionales independientes y los establecimientos comerciales tienen poca visibilidad para darse a conocer y atraer nuevos clientes fuera de su c\'irculo cercano. \\
\hline
\textbf{Oportunidad} & Crear una plataforma organizada donde los profesionales independientes y los locales comerciales puedan publicar un perfil digital con sus servicios, fotograf\'ias de trabajos previos, ubicaci\'on por sector, horarios y canales de contacto directo, facilitando que los clientes encuentren opciones cercanas y confiables en poco tiempo. \\
\hline
\textbf{Objetivo general} & Identificar las necesidades, dificultades, expectativas y caracter\'isticas de b\'usqueda y promoci\'on de servicios de los clientes, profesionales independientes y propietarios de establecimientos comerciales, con la finalidad de establecer los requisitos del sistema propuesto. \\
\hline
\end{tabularx}
\end{table}

% ==================== A.1 VISION Y OBJETIVOS DE NEGOCIO ====================
\subsection*{A.1. Declaraci\'on de visi\'on}
\addcontentsline{toc}{subsection}{A.1. Declaraci\'on de visi\'on}

Para los vecinos de Quevedo que hoy tienen que buscar un electricista, un plomero o un local preguntando a conocidos o revisando grupos de Facebook, NEXORA ser\'a un directorio local donde cada profesional y cada negocio muestra sus servicios, fotos de trabajos anteriores, sector y n\'umero de contacto. As\'i, el cliente puede encontrar a alguien de confianza cerca de su casa en pocos minutos, y el proveedor consigue clientes fuera de su c\'irculo de siempre.

\subsection*{A.2. Objetivos de negocio}
\addcontentsline{toc}{subsection}{A.2. Objetivos de negocio}

\begin{enumerate}[label=\textbf{OB-\arabic*.}, leftmargin=3.2em, itemsep=4pt]
    \item \textbf{Reducir el tiempo de b\'usqueda.} Que un cliente encuentre y contacte a un profesional o local en menos de 10 minutos. Hoy esto puede tomar uno o m\'as d\'ias (es un supuesto que se confirmar\'a con las entrevistas). La meta es lograrlo dentro de los primeros 6 meses despu\'es del lanzamiento. Se medir\'a con el tiempo que pasa desde que el cliente hace la b\'usqueda hasta que se comunica con el proveedor, por WhatsApp o por llamada.
    \item \textbf{Tener proveedores registrados.} Lograr al menos 150 perfiles completos y publicados (100 profesionales independientes y 50 establecimientos comerciales) en los primeros 4 meses de funcionamiento. El equipo estim\'o esa cantidad para cubrir los oficios m\'as pedidos en los principales sectores de Quevedo, y se ajustar\'a con la elicitaci\'on. Solo cuentan los perfiles que tengan servicios, sector y medio de contacto.
    \item \textbf{Ofrecer opciones cercanas.} Conseguir que, al final del primer a\~no, el 70\,\% de las b\'usquedas hechas en la plataforma muestren al menos un proveedor del mismo sector del cliente o de un sector colindante (uno que comparte l\'imite con el suyo). Se medir\'a con los registros de b\'usqueda filtrados por ubicaci\'on.
\end{enumerate}

% ==================== B. DESCRIPCION DE LAS PARTES INTERESADAS ====================
\section*{B. Descripci\'on de las partes interesadas}
\addcontentsline{toc}{section}{B. Descripci\'on de las partes interesadas}

De acuerdo con la norma ISO/IEC/IEEE 29148:2018 \cite{iso29148} y el programa de estudio del IREB \cite{ireb2026}, identificar a las partes interesadas permite conocer qui\'enes utilizar\'an el sistema o se ver\'an afectados por su funcionamiento. En este proyecto se identificaron ocho partes interesadas: tres grupos de usuarios externos (clientes, profesionales independientes y propietarios de establecimientos comerciales), los roles encargados de la administraci\'on y del an\'alisis del sistema, la organizaci\'on cliente que acept\'o participar en el levantamiento de requisitos, el docente de la asignatura y la autoridad de protecci\'on de datos personales. En la Tabla~\ref{tab:stakeholders} se describen sus intereses, su nivel de influencia y la estrategia para involucrarlas en el proceso.

\begin{table}[H]
\centering
\caption{Descripci\'on y an\'alisis de las partes interesadas del sistema}
\label{tab:stakeholders}
\footnotesize
\renewcommand{\arraystretch}{1.1}
\begin{tabularx}{\textwidth}{|L{3.1cm}|Y|C{2.0cm}|Y|}
\hline
\rowcolor{tableheader}
\textbf{Rol} & \textbf{Intereses principales en el sistema} & \textbf{Nivel de influencia} & \textbf{Estrategia de involucramiento} \\
\hline
\textbf{1. Clientes o usuarios} & Encontrar con facilidad profesionales o establecimientos comerciales dentro de su ciudad, cant\'on o sector; consultar servicios, horarios, ubicaci\'on y fotograf\'ias de trabajos realizados que les generen confianza antes de comunicarse. & Media & Aplicar entrevistas y encuestas para conocer c\'omo buscan servicios actualmente, qu\'e dificultades tienen y qu\'e informaci\'on necesitan ver antes de contactar a un proveedor. \\
\hline
\textbf{2. Profesionales independientes} & Dar a conocer su oficio o actividad (electricidad, plomer\'ia, cerrajer\'ia, carpinter\'ia, entre otros), mostrar fotograf\'ias de trabajos realizados, indicar los sectores a los que pueden desplazarse y conseguir nuevos clientes. & Media & Realizar entrevistas individuales para identificar c\'omo promocionan hoy su trabajo, qu\'e problemas enfrentan para captar clientes y qu\'e datos desean incluir en su perfil digital. \\
\hline
\textbf{3. Propietarios de establecimientos comerciales} & Promocionar su local f\'isico, difundir los servicios que ofrecen, mostrar su direcci\'on, horarios de atenci\'on, fotograf\'ias del establecimiento y n\'umeros de contacto para atraer a personas que a\'un no conocen su negocio. & Media & Realizar entrevistas en sus locales comerciales para conocer qu\'e medios digitales utilizan, qu\'e limitaciones tienen para dar visibilidad a su negocio y qu\'e caracter\'isticas requieren en su perfil. \\
\hline
\textbf{4. Administrador del sistema} & Mantener organizadas las categor\'ias de servicios y los sectores geogr\'aficos dentro de la plataforma, as\'i como revisar que los perfiles publicados contengan informaci\'on clara y real. & Alta & Definir junto con el equipo las reglas b\'asicas para el registro, organizaci\'on y publicaci\'on de perfiles dentro de la plataforma. \\
\hline
\textbf{5. Equipo de an\'alisis de requisitos} & Recopilar informaci\'on real de los usuarios mediante los instrumentos de elicitaci\'on, analizar sus necesidades y documentar de forma clara los requisitos del sistema propuesto. & Alta & Planificar y ejecutar las sesiones de entrevistas y encuestas, organizar los hallazgos obtenidos y validar que el alcance responda a lo expresado por los participantes. \\
\hline
\textbf{6. Organizaci\'on cliente} (empresa de sistemas de pesaje, Quevedo) & Que la plataforma ayude a que m\'as personas conozcan el negocio, ya que actualmente los clientes llegan casi solo por recomendaci\'on. & Alta & Como firm\'o la carta de aceptaci\'on, nos reuniremos con el propietario para mostrarle los avances y confirmar si el proyecto va por buen camino. Adem\'as, sigue participando en las entrevistas. \\
\hline
\textbf{7. Docente de la asignatura} & Que el proyecto cumpla con lo que exige la asignatura y que los entregables se presenten completos, bien elaborados y a tiempo. & Alta & Presentar cada entregable en las fechas establecidas por la asignatura y aplicar las correcciones que indique. \\
\hline
\textbf{8. Superintendencia de Protecci\'on de Datos Personales} & Que los datos personales que aparezcan en los perfiles, como nombres, tel\'efonos y direcciones, se publiquen conforme a la LOPDP~\cite{lopdp2021}. & Alta & No existe contacto directo. El equipo revisa la LOPDP~\cite{lopdp2021} y la toma en cuenta al definir qu\'e datos se muestran en cada perfil. \\
\hline
\end{tabularx}
\end{table}

\subsection*{B.1. Poder, inter\'es y disponibilidad}
\addcontentsline{toc}{subsection}{B.1. Poder, inter\'es y disponibilidad}

Para decidir c\'omo gestionar a cada parte interesada se valor\'o su poder sobre el proyecto y su inter\'es en el sistema, y se registr\'o cu\'ando y por qu\'e medio se le puede contactar (Tabla~\ref{tab:poder-interes}). Con esa valoraci\'on se construy\'o la matriz poder--inter\'es de la Figura~\ref{fig:poder-interes}.

\begin{table}[H]
\centering
\caption{Poder, inter\'es y disponibilidad de las partes interesadas}
\label{tab:poder-interes}
\small
\begin{tabularx}{\textwidth}{|L{3.4cm}|C{1.3cm}|C{1.3cm}|L{2.4cm}|Y|}
\hline
\rowcolor{tableheader}
\textbf{Parte interesada} & \textbf{Poder} & \textbf{Inter\'es} & \textbf{Cuadrante} & \textbf{Disponibilidad} \\
\hline
1. Clientes o usuarios & Bajo & Alto & Mantener informados & En las tardes y fines de semana, por encuesta en l\'inea o en persona, como se hizo con Nexora-01 y Nexora-03. \\
\hline
2. Profesionales independientes & Bajo & Alto & Mantener informados & Durante el d\'ia suelen estar trabajando; es mejor escribirles por WhatsApp o llamarlos en la noche, como se hizo con Nexora-02. \\
\hline
3. Propietarios de establecimientos comerciales & Bajo & Alto & Mantener informados & Visita al local en horario de atenci\'on, en momentos con poca afluencia de clientes. \\
\hline
4. Administrador del sistema & Alto & Alto & Gestionar de cerca & Rol a\'un no asignado; mientras tanto lo asume el equipo del proyecto, que define las reglas de categor\'ias y perfiles. Debe consultarse antes de cerrar esos requisitos. \\
\hline
5. Equipo de an\'alisis de requisitos & Alto & Alto & Gestionar de cerca & Disponible de forma permanente; coordinaci\'on en clase, por WhatsApp o Google Meet. \\
\hline
6. Organizaci\'on cliente & Alto & Alto & Gestionar de cerca & El propietario atiende en la tarde, despu\'es de la jornada de trabajo, en persona o por llamada. \\
\hline
7. Docente de la asignatura & Alto & Bajo & Mantener satisfechos & Inter\'es bajo en usar la plataforma, porque no es usuario del sistema; su inter\'es est\'a en el proceso y en la calidad de los entregables. Disponible en horas de clase y tutor\'ias, o por correo institucional. \\
\hline
8. Superintendencia de Protecci\'on de Datos Personales & Alto & Bajo & Mantener satisfechos & Sin contacto directo; lo necesario est\'a en la LOPDP~\cite{lopdp2021} y en su sitio web, que se puede consultar en cualquier momento. \\
\hline
\end{tabularx}
\end{table}

\begin{figure}[H]
\centering
\resizebox{0.86\textwidth}{!}{% Matriz poder-interes de las partes interesadas (TA-PFC-E1).
% Fuente editable de la Figura de la seccion B. Se incluye en main.tex con \input.
% Numeracion de las partes interesadas segun la Tabla de la seccion B.
\begin{tikzpicture}[
  quad/.style={draw=black!60, minimum width=6.4cm, minimum height=3.6cm, align=center},
  sh/.style={circle, draw=uteqgreen, fill=white, very thick, minimum size=0.75cm, inner sep=0pt, font=\small\bfseries}
]
% Cuadrantes
\node[quad, fill=uteqgreen!8]  (q2) at (0,3.6)   {};
\node[quad, fill=uteqgreen!22] (q1) at (6.4,3.6) {};
\node[quad, fill=black!3]      (q4) at (0,0)     {};
\node[quad, fill=uteqgreen!8]  (q3) at (6.4,0)   {};
% Titulos de cuadrante
\node[font=\small\bfseries, anchor=north] at ([yshift=-4pt]q2.north) {Mantener satisfechos};
\node[font=\small\bfseries, anchor=north] at ([yshift=-4pt]q1.north) {Gestionar de cerca};
\node[font=\small\bfseries, anchor=north] at ([yshift=-4pt]q4.north) {Monitorear};
\node[font=\small\bfseries, anchor=north] at ([yshift=-4pt]q3.north) {Mantener informados};
% Partes interesadas
\node[sh] at (-0.8,3.3) {7};
\node[sh] at ( 0.8,3.3) {8};
\node[sh] at (5.4,3.3) {4};
\node[sh] at (6.4,3.3) {5};
\node[sh] at (7.4,3.3) {6};
\node[sh] at (5.4,-0.3) {1};
\node[sh] at (6.4,-0.3) {2};
\node[sh] at (7.4,-0.3) {3};
\node[font=\footnotesize\itshape, text=black!55] at (0,-0.3) {(ninguna)};
% Ejes
\draw[-latex, thick] (-3.4,-1.95) -- (9.9,-1.95) node[midway, below=12pt, font=\small\bfseries] {Inter\'es en el sistema};
\draw[-latex, thick] (-3.4,-1.95) -- (-3.4,5.55);
\node[rotate=90, font=\small\bfseries] at (-4.75,1.8) {Poder sobre el proyecto};
\node[font=\footnotesize, anchor=north] at (0,-1.95) {Bajo};
\node[font=\footnotesize, anchor=north] at (6.4,-1.95) {Alto};
\node[font=\footnotesize, anchor=east] at (-3.5,0) {Bajo};
\node[font=\footnotesize, anchor=east] at (-3.5,3.6) {Alto};
\end{tikzpicture}
}
\caption{Matriz poder--inter\'es de las partes interesadas (numeraci\'on de la Tabla~\ref{tab:stakeholders})}
\label{fig:poder-interes}
\end{figure}

% ==================== C. DEFINICION DEL ALCANCE ====================
\section*{C. Definici\'on del alcance}
\addcontentsline{toc}{section}{C. Definici\'on del alcance}

Para evitar confusiones sobre el prop\'osito del sistema, es necesario definir con claridad hasta d\'onde llega su funcionamiento \cite{wiegers2013}. El sistema propuesto tiene como l\'imite ser una plataforma informativa de \textbf{promoci\'on, difusi\'on, b\'usqueda y contacto directo} entre clientes y proveedores de servicios locales. Esto significa que el sistema ayuda al cliente a encontrar, comparar y elegir a un profesional independiente o establecimiento comercial mediante su perfil digital, pero la comunicaci\'on final, el acuerdo del trabajo y el pago se realizan de forma directa entre ambas partes fuera de la plataforma.

En la Tabla~\ref{tab:alcance} se detalla lo que est\'a dentro del alcance del sistema y lo que queda fuera de \'el.

\begin{table}[H]
\centering
\caption{Delimitaci\'on de lo que est\'a dentro y fuera del alcance del sistema}
\label{tab:alcance}
\small
\begin{tabularx}{\textwidth}{|Y|Y|}
\hline
\rowcolor{tableheader}
\textbf{Dentro del alcance (Lo que el sistema S\'I har\'a)} & \textbf{Fuera del alcance (Lo que el sistema NO har\'a)} \\
\hline
$\bullet$ \textbf{Perfiles para profesionales independientes:} Registro y publicaci\'on de un perfil digital con datos del trabajador, oficio, servicios ofrecidos, experiencia, sectores a los que se desplaza y medios de contacto. & $\bullet$ \textbf{Cobros o pagos dentro de la plataforma:} El sistema no procesar\'a pagos con tarjetas de cr\'edito, transferencias bancarias ni cobrar\'a comisiones por los servicios. \\
\hline
$\bullet$ \textbf{Perfiles para establecimientos comerciales:} Registro y publicaci\'on de un perfil para negocios con nombre del local, servicios que ofrece, direcci\'on f\'isica, ubicaci\'on, horarios de atenci\'on y n\'umeros de contacto. & $\bullet$ \textbf{Negociaci\'on de precios o contrataci\'on interna:} El sistema no calcular\'a presupuestos ni generar\'a contratos; el cliente y el proveedor acuerdan el servicio directamente al comunicarse. \\
\hline
$\bullet$ \textbf{Galer\'ia de fotograf\'ias de trabajos realizados:} Espacio en el perfil para subir y mostrar fotograf\'ias reales de trabajos anteriores o de las instalaciones del establecimiento comercial. & $\bullet$ \textbf{Rastreo en vivo del trabajador:} No se realizar\'a seguimiento por GPS en tiempo real del desplazamiento del profesional hacia la casa del cliente. \\
\hline
$\bullet$ \textbf{B\'usqueda y filtrado de servicios:} Buscador para que los clientes encuentren opciones filtrando por tipo de servicio u oficio y por ubicaci\'on geogr\'afica (ciudad, cant\'on o sector). & $\bullet$ \textbf{Facturaci\'on electr\'onica o contabilidad:} La plataforma no emitir\'a facturas ni llevar\'a el control contable de los ingresos de los profesionales o negocios. \\
\hline
$\bullet$ \textbf{Informaci\'on de confianza y contacto directo:} Visualizaci\'on clara de los datos del proveedor y enlaces directos para comunicarse por WhatsApp, llamada telef\'onica o visitar el local. & $\bullet$ \textbf{Intermediaci\'on legal o garant\'ias de obra:} El sistema funciona como un medio de difusi\'on y directorio, sin intervenir legalmente en la ejecuci\'on f\'isica del trabajo. \\
\hline
\end{tabularx}
\end{table}

\textbf{Restricciones principales del sistema:}
\begin{itemize}[leftmargin=1.5em, itemsep=2pt]
    \item \textbf{Privacidad y autorizaci\'on de datos:} Toda informaci\'on publicada en los perfiles (n\'umeros de contacto, direcci\'on, servicios y fotograf\'ias de trabajos realizados) ser\'a registrada voluntariamente por cada profesional o propietario de establecimiento para fines de promoci\'on p\'ublica.
    \item \textbf{Veracidad y actualizaci\'on de la informaci\'on:} Los perfiles requerir\'an completar los datos esenciales (servicios, sector o direcci\'on y medio de contacto) antes de mostrarse al p\'ublico para evitar anuncios vac\'ios o desactualizados.
    \item \textbf{Tiempo de ejecuci\'on acad\'emica:} El levantamiento, an\'alisis y especificaci\'on de requisitos debe ajustarse a los plazos establecidos dentro del per\'iodo acad\'emico de la asignatura.
\end{itemize}

% ==================== D. DIAGRAMA DE CONTEXTO DEL SISTEMA ====================
\section*{D. Diagrama de contexto del sistema}
\addcontentsline{toc}{section}{D. Diagrama de contexto del sistema}

El diagrama de contexto permite visualizar los l\'imites de la plataforma propuesta, ubic\'andola en el centro como un \'unico elemento y mostrando qu\'e informaci\'on intercambia con cada uno de los actores externos que interact\'uan con ella \cite{ireb2026}. En la Figura~\ref{fig:contexto} se presenta el diagrama de contexto del sistema.

\begin{figure}[H]
\centering
\resizebox{\textwidth}{!}{%
\begin{tikzpicture}[
    actorbox/.style={rectangle, draw=black!85, fill=white, thick, minimum width=3.2cm, minimum height=1.2cm, align=center, font=\footnotesize\bfseries},
    sysnode/.style={circle, draw=uteqgreen!90!black, fill=tableheader, very thick, minimum size=3.4cm, align=center, font=\footnotesize\bfseries},
    flow/.style={-{Stealth[length=2.2mm]}, semithick, draw=black!85},
    lbl/.style={font=\scriptsize, align=center, text width=4.5cm}
]

% Nodo central: Sistema
\node[sysnode] (sys) at (0,0) {Plataforma de\\Promoci\'on y Difusi\'on\\de Servicios\\(NEXORA)};

% Actores en las 4 esquinas
\node[actorbox] (cli) at (-6.7, 2.75) {Cliente o\\Usuario};
\node[actorbox] (pro) at (6.7, 2.75) {Profesional\\Independiente};
\node[actorbox] (com) at (-6.7, -2.75) {Propietario de\\Establecimiento\\Comercial};
\node[actorbox] (adm) at (6.7, -2.75) {Administrador\\del Sistema};

% 1. Cliente <-> Sistema (Arriba Izquierda)
\draw[flow] ([yshift=0.3cm]cli.east) -- (-1.05, 3.05) -- (-1.05, 1.54);
\node[lbl, text width=4.1cm, anchor=south] at (-3.1, 3.12) {B\'usqueda de servicio u oficio, filtro por ciudad, cant\'on o sector y consulta de perfil.};

\draw[flow] (-1.28, 1.36) -- (-6.7, 1.36) -- (cli.south);
\node[lbl, text width=4.5cm, anchor=north] at (-4.3, 1.30) {Perfiles de profesionales y locales, fotos de trabajos realizados, ubicaci\'on, horarios y enlaces de contacto directo (WhatsApp o llamada).};

% 2. Profesional Independiente <-> Sistema (Arriba Derecha)
\draw[flow] ([yshift=0.3cm]pro.west) -- (1.05, 3.05) -- (1.05, 1.54);
\node[lbl, text width=4.1cm, anchor=south] at (3.1, 3.12) {Datos de oficio, servicios ofrecidos, experiencia, fotos de trabajos previos, zonas y contacto.};

\draw[flow] (1.28, 1.36) -- (6.7, 1.36) -- (pro.south);
\node[lbl, text width=4.5cm, anchor=north] at (4.3, 1.30) {Perfil digital publicado en el directorio, sectores visibles y confirmaci\'on de registro.};

% 3. Propietario de Establecimiento <-> Sistema (Abajo Izquierda)
\draw[flow] (com.north) -- (-6.7, -1.36) -- (-1.28, -1.36);
\node[lbl, text width=4.5cm, anchor=south] at (-4.3, -1.30) {Datos del establecimiento, servicios, direcci\'on, horarios, fotograf\'ias y medios de contacto.};

\draw[flow] (-1.05, -1.54) -- (-1.05, -3.05) -- ([yshift=-0.3cm]com.east);
\node[lbl, text width=4.1cm, anchor=north] at (-3.1, -3.12) {Perfil comercial publicado, ubicaci\'on visible por sector y confirmaci\'on de registro.};

% 4. Administrador <-> Sistema (Abajo Derecha)
\draw[flow] (adm.north) -- (6.7, -1.36) -- (1.28, -1.36);
\node[lbl, text width=4.5cm, anchor=south] at (4.3, -1.30) {Organizaci\'on de categor\'ias de servicios, sectores geogr\'aficos y revisi\'on de perfiles.};

\draw[flow] (1.05, -1.54) -- (1.05, -3.05) -- ([yshift=-0.3cm]adm.west);
\node[lbl, text width=4.1cm, anchor=north] at (3.1, -3.12) {Listado de perfiles registrados y estado de las publicaciones en la plataforma.};

\end{tikzpicture}%
}
\caption{Diagrama de contexto del sistema propuesto}
\label{fig:contexto}
\end{figure}

Para complementar el diagrama de contexto, en la Tabla~\ref{tab:perspectivas} se sintetiza el an\'alisis del entorno de NEXORA mediante las ocho perspectivas de contexto propuestas por Pohl~\cite{pohl2025}. Cada fila registra la parte interesada consultada, el hallazgo respaldado en las entrevistas o documentos del proyecto y el requisito preliminar derivado; los aspectos que requieren evidencia adicional o acuerdo del equipo se se\~nalan como pendientes de confirmar.

\begin{table}[H]
\centering
\caption{An\'alisis del contexto del sistema seg\'un las ocho perspectivas de Pohl}
\label{tab:perspectivas}
\footnotesize
\renewcommand{\arraystretch}{1.15}
\begin{tabularx}{\textwidth}{|L{2.05cm}|L{2.55cm}|>{\hsize=1.12\hsize}Y|>{\hsize=0.88\hsize}Y|}
\hline
\rowcolor{tableheader}
\textbf{Perspectiva} & \textbf{Parte interesada a entrevistar} & \textbf{Hallazgos clave} & \textbf{Requisitos derivados (preliminares)} \\
\hline
\textbf{Uso} & Cliente o usuario & El cliente entrevistado indica que las fotograf\'ias de trabajos anteriores le permiten valorar la calidad del servicio (Entrevista al cliente, pregunta 5). & La plataforma deber\'a permitir que el cliente consulte fotograf\'ias de trabajos anteriores en el perfil del profesional. \\
\hline
\textbf{Sujeto} & Profesional independiente & El profesional independiente (herrero) ofrece fabricaci\'on de estructuras met\'alicas, soldadura y reparaciones; atiende en Quevedo y sectores cercanos, y decide hasta d\'onde desplazarse seg\'un el tipo de trabajo y la distancia (Entrevista al profesional independiente, preguntas 8 y 13). & La plataforma deber\'a permitir que el profesional registre sus servicios y los sectores donde atiende. \\
\hline
\textbf{TI} & Cliente o usuario y profesional independiente & El cliente utiliza WhatsApp para buscar servicios (Entrevista al cliente, pregunta 3) y el profesional independiente lo emplea para enviar fotograf\'ias a quienes le consultan (Entrevista al profesional independiente, pregunta 11). Seg\'un el alcance (secci\'on C, Tabla~\ref{tab:alcance}), la comunicaci\'on final ocurre fuera de la plataforma mediante enlaces directos. & El perfil deber\'a ofrecer un enlace de contacto externo por WhatsApp cuando el proveedor registre ese canal, conforme al alcance de la secci\'on C. \\
\hline
\textbf{Desarrollo} & Equipo del proyecto & El equipo mantiene el informe en \LaTeX{} en el repositorio y document\'o c\'omo compilarlo (\texttt{README.md} y \texttt{PFC/E1/main.tex}), sujeto a los plazos acad\'emicos (secci\'on C). \textbf{Pendiente de confirmar:} acuerdo del equipo sobre metodolog\'ia de desarrollo, herramientas y proceso de validaci\'on. & El equipo deber\'a poder regenerar el PDF del informe siguiendo las instrucciones de \texttt{README.md}. \textbf{Pendiente de confirmar:} requisitos de proceso sobre metodolog\'ia, herramientas y validaci\'on. \\
\hline
\textbf{Negocio} & Profesional independiente & El profesional independiente (herrero) consigue clientes principalmente por recomendaciones y tiene dificultad para darse a conocer en otros sectores (Entrevista al profesional independiente, preguntas 9 y 10). & La plataforma deber\'a permitir buscar proveedores por tipo de servicio y sector para que los clientes descubran profesionales que atienden en su zona. \\
\hline
\textbf{Innovaci\'on} & Profesional independiente & Ante la idea de un perfil digital, el profesional independiente (herrero) expres\'o que le gustar\'ia reunir fotograf\'ias de trabajos terminados y una lista clara de lo que puede hacer (Entrevista al profesional independiente, pregunta 14). & La plataforma deber\'a reunir en un mismo perfil p\'ublico la lista de servicios, las fotograf\'ias de trabajos y el medio de contacto que el profesional elija. \\
\hline
\textbf{Sociedad} & Profesional independiente & El profesional independiente (herrero) expres\'o inter\'es en incluir su tel\'efono y direcci\'on en un perfil digital (Entrevista al profesional independiente, pregunta 14), acorde con el registro voluntario de la secci\'on C. \textbf{Pendiente de confirmar:} qu\'e datos autorizar\'ia a publicar en un perfil real y obligaciones aplicables de la LOPDP~\cite{lopdp2021}. & La plataforma deber\'a mostrar \'unicamente los datos de contacto que el proveedor decida publicar. \textbf{Pendiente de confirmar:} mecanismo de autorizaci\'on expresa, retiro de datos seg\'un la LOPDP~\cite{lopdp2021} y accesibilidad. \\
\hline
\textbf{T\'ecnica} & Profesional independiente & El profesional independiente (herrero) guarda en su celular las fotograf\'ias que comparte con sus clientes (Entrevista al profesional independiente, pregunta 11). \textbf{Pendiente de confirmar:} limitaciones de dispositivo, navegador o conectividad en Quevedo y sectores cercanos. & La plataforma deber\'a permitir que el profesional cargue fotograf\'ias a su perfil desde el dispositivo que utilice. \textbf{Pendiente de confirmar:} restricciones t\'ecnicas de tama\~no de im\'agenes y conectividad limitada. \\
\hline
\end{tabularx}
\end{table}

% ==================== E. PLAN DE ELICITACION ====================
\section*{E. Plan de elicitaci\'on}
\addcontentsline{toc}{section}{E. Plan de elicitaci\'on}

Para levantar los requisitos del sistema a partir de la experiencia real de las personas involucradas, se estructur\'o un plan de elicitaci\'on que combina entrevistas semiestructuradas, observaci\'on no participante y encuestas \cite{wiegers2013, pohl2025}.

\subsection*{1. T\'ecnicas seleccionadas y justificaci\'on}
\begin{itemize}[leftmargin=1.5em, itemsep=3pt]
    \item \textbf{Entrevistas semiestructuradas:} Se seleccion\'o esta t\'ecnica porque permite conversar directamente con los tres grupos principales de partes interesadas (clientes, profesionales independientes y propietarios de establecimientos comerciales). Mediante un guion de 20 preguntas divididas por perfil, se puede profundizar en c\'omo buscan o promocionan servicios actualmente, qu\'e dificultades enfrentan, qu\'e valor le dan a las fotograf\'ias de trabajos anteriores y a la ubicaci\'on por sector, y qu\'e informaci\'on esperan encontrar o mostrar en un perfil digital.
    \item \textbf{Observaci\'on no participante:} Se incorpor\'o para observar en el entorno real c\'omo el propietario o encargado de la organizaci\'on cliente muestra sus servicios, atiende una consulta y comunica c\'omo contactarlo (o, si se aplica a un cliente, c\'omo busca y eval\'ua un proveedor antes de comunicarse), sin intervenir en la actividad. Esta t\'ecnica complementa las entrevistas porque permite identificar conocimiento t\'acito, herramientas reales y dificultades que los participantes no siempre verbalizan \cite{pohl2025, ireb2026}. Su ficha de registro corresponde a \texttt{PFC/E1/instrumentos/guia-observacion.tex}.
    \item \textbf{Encuestas / Cuestionarios:} Se eligi\'o como t\'ecnica complementaria para recoger la opini\'on de un n\'umero mayor de ciudadanos y usuarios del cant\'on Quevedo en poco tiempo. Permite confirmar con datos claros cu\'ales son los servicios m\'as buscados, los medios de comunicaci\'on preferidos y los aspectos que generan mayor confianza antes de contactar a un profesional o negocio. El cuestionario tiene 12 preguntas en cinco bloques (datos generales, c\'omo se busca hoy, dificultades, confianza e informaci\'on, y cierre); cada pregunta indica qu\'e hallazgo de las entrevistas Nexora-01 y Nexora-03 comprueba. Es an\'onimo, voluntario y solo para mayores de 18 a\~nos, y corresponde a \texttt{PFC/E1/\allowbreak instrumentos/\allowbreak cuestionario.tex}.
\end{itemize}

\subsection*{2. Participantes, fechas tentativas y artefactos esperados}

En la Tabla~\ref{tab:plan_elicitacion} se detalla la planificaci\'on de las actividades de elicitaci\'on, indicando los participantes de cada sesi\'on, las fechas tentativas de ejecuci\'on y los artefactos que se obtendr\'an como resultado.

\begin{table}[H]
\centering
\caption{Planificaci\'on de actividades de elicitaci\'on de requisitos}
\label{tab:plan_elicitacion}
\footnotesize
\renewcommand{\arraystretch}{1.12}
\begin{tabularx}{\textwidth}{|L{2.7cm}|L{3.1cm}|Y|C{2.2cm}|L{3.2cm}|}
\hline
\rowcolor{tableheader}
\textbf{T\'ecnica} & \textbf{Participantes} & \textbf{Objetivo de la actividad} & \textbf{Fechas tentativas} & \textbf{Artefactos esperados} \\
\hline
\textbf{Entrevista semiestructurada (Bloque A)} & \textbf{Parte interesada 1:} Clientes o usuarios (Preguntas 1 a 7) & Conocer c\'omo buscan actualmente profesionales o locales comerciales, qu\'e dificultades tienen, la importancia de las fotos de trabajos previos, la ubicaci\'on y qu\'e les genera confianza. & 20 de noviembre de 2026 & Guion de entrevista aplicado, registro de respuestas de clientes y necesidades de b\'usqueda identificadas. \\
\hline
\textbf{Entrevista semiestructurada (Bloque B)} & \textbf{Parte interesada 2:} Profesionales independientes (Preguntas 8 a 14) & Conocer a qu\'e oficio se dedican, c\'omo consiguen clientes, qu\'e medios usan para mostrar fotos de sus trabajos, a qu\'e sectores se desplazan y qu\'e desean en su perfil digital. & 20 de noviembre de 2026 & Guion de entrevista aplicado, registro de respuestas de profesionales y requisitos para el perfil de oficios. \\
\hline
\textbf{Entrevista semiestructurada (Bloque C)} & \textbf{Parte interesada 3:} Propietarios de establecimientos comerciales (Preguntas 15 a 20) & Identificar qu\'e servicios ofrecen, c\'omo promocionan su local, qu\'e dificultades tienen para darse a conocer y qu\'e datos (direcci\'on, horarios, fotos, contacto) requieren mostrar. & 20 de noviembre de 2026 & Guion de entrevista aplicado, registro de respuestas de comerciantes y requisitos para el perfil comercial. \\
\hline
\textbf{Observaci\'on no participante} & \textbf{Parte interesada 6 (y 3):} Propietario o encargado de la organizaci\'on cliente (empresa de servicios t\'ecnicos de balanzas y b\'asculas); opcionalmente clientes o usuarios (Parte interesada 1) & Observar sin intervenir c\'omo el proveedor muestra sus servicios, atiende una consulta y comunica c\'omo contactarlo (o c\'omo un cliente busca y revisa datos antes de contactar), identificando canales, informaci\'on disponible y dificultades visibles. & Por confirmar & Gu\'ia de observaci\'on aplicada (\texttt{guia-\allowbreak observacion.tex}), registro separado de hechos observados e interpretaciones y notas anonimizadas (m\'as evidencia audiovisual solo si se autoriza). \\
\hline
\textbf{Encuesta a usuarios} & Ciudadanos y usuarios del cant\'on Quevedo & Recopilar preferencias generales sobre frecuencia de b\'usqueda de servicios, canales de contacto preferidos y filtros de ubicaci\'on. & 23 al 27 de noviembre de 2026 & Cuestionario aplicado, resumen de resultados y listado preliminar de requisitos del sistema. \\
\hline
\end{tabularx}
\end{table}

\subsection*{3. Guion de entrevista dise\~nado por parte interesada}

\textbf{Objetivo:} Identificar las necesidades, dificultades, expectativas y caracter\'isticas de b\'usqueda y promoci\'on de servicios de los clientes, profesionales independientes y propietarios de establecimientos comerciales, con la finalidad de establecer los requisitos del sistema propuesto.

\vspace{3pt}
\noindent\textbf{A. Parte interesada 1: Clientes o usuarios}\\
\textit{Objetivo de la entrevista:} Conocer c\'omo las personas buscan actualmente profesionales o establecimientos comerciales cuando necesitan un servicio.
\begin{enumerate}[label=\textbf{\arabic*.}, leftmargin=1.8em, itemsep=1pt, topsep=2pt]
    \item Cuando necesita contratar un servicio, como electricidad, plomer\'ia, cerrajer\'ia o carpinter\'ia, ?`c\'omo suele encontrar a la persona o establecimiento que lo ofrece?
    \item ?`Qu\'e dificultades ha experimentado al momento de buscar un profesional o establecimiento que ofrezca el servicio que necesita?
    \item ?`Qu\'e medios utiliza actualmente para encontrar profesionales o negocios? Por ejemplo, recomendaciones de conocidos, redes sociales, buscadores de Internet u otros.
    \item ?`Qu\'e informaci\'on considera importante conocer antes de comunicarse con un profesional o establecimiento comercial?
    \item ?`Qu\'e importancia tienen para usted las fotograf\'ias de trabajos realizados anteriormente al momento de elegir un profesional? ?`Por qu\'e?
    \item ?`C\'omo influye la ubicaci\'on geogr\'afica al momento de buscar un servicio? ?`Qu\'e tan importante es encontrar profesionales dentro de su ciudad, cant\'on o sector?
    \item ?`Qu\'e aspectos le generan confianza o desconfianza cuando encuentra a un profesional o establecimiento que no conoce?
\end{enumerate}

\vspace{3pt}
\noindent\textbf{B. Parte interesada 2: Profesionales independientes}\\
\textit{Objetivo de la entrevista:} Conocer c\'omo los trabajadores independientes ofrecen sus servicios, promocionan sus habilidades y consiguen nuevos clientes.
\begin{enumerate}[label=\textbf{\arabic*.}, leftmargin=1.8em, itemsep=1pt, topsep=2pt, start=8]
    \item ?`Podr\'ia contarme a qu\'e actividad u oficio se dedica y qu\'e tipos de servicios ofrece actualmente?
    \item ?`C\'omo consigue actualmente nuevos clientes y c\'omo da a conocer los servicios que ofrece?
    \item ?`Qu\'e dificultades encuentra al momento de promocionar sus servicios o conseguir nuevos clientes?
    \item ?`Qu\'e medios utiliza actualmente para mostrar fotograf\'ias, ejemplos o evidencias de los trabajos que ha realizado?
    \item ?`Qu\'e informaci\'on considera importante que un cliente conozca sobre usted antes de comunicarse para solicitar sus servicios?
    \item ?`En qu\'e lugares o sectores ofrece sus servicios? ?`C\'omo determina las zonas a las que est\'a dispuesto a desplazarse para realizar un trabajo?
    \item Si tuviera un perfil digital donde pudiera mostrar sus servicios, experiencia, fotograf\'ias de trabajos realizados y datos de contacto, ?`qu\'e informaci\'on le gustar\'ia incluir y qu\'e beneficios esperar\'ia obtener?
\end{enumerate}

\vspace{3pt}
\noindent\textbf{C. Parte interesada 3: Propietarios de establecimientos comerciales}\\
\textit{Objetivo de la entrevista:} Identificar c\'omo los negocios promocionan sus servicios, dan a conocer sus establecimientos y atraen nuevos clientes.
\begin{enumerate}[label=\textbf{\arabic*.}, leftmargin=1.8em, itemsep=1pt, topsep=2pt, start=15]
    \item ?`Podr\'ia describirme qu\'e servicios ofrece su establecimiento y qu\'e tipo de clientes suelen acudir a su negocio?
    \item ?`C\'omo promociona actualmente su establecimiento y c\'omo consigue que nuevos clientes conozcan los servicios que ofrece?
    \item ?`Qu\'e dificultades tiene para dar visibilidad a su negocio y llegar a personas que todav\'ia no conocen su establecimiento?
    \item ?`Qu\'e informaci\'on considera importante que las personas puedan consultar antes de visitar o comunicarse con su establecimiento? Por ejemplo, direcci\'on, horarios, servicios, fotograf\'ias o n\'umeros de contacto.
    \item ?`Utiliza actualmente medios digitales para promocionar su negocio? Si es as\'i, ?`qu\'e herramientas utiliza y qu\'e dificultades ha encontrado al utilizarlas?
    \item Si pudiera contar con un perfil digital para mostrar su establecimiento, los servicios que ofrece, fotograf\'ias, ubicaci\'on y medios de contacto, ?`qu\'e caracter\'isticas considerar\'ia necesarias?
\end{enumerate}

% ==================== F. PROCEDIMIENTO DE ANONIMIZACION Y CUSTODIA ====================
\section*{F. Procedimiento de anonimizaci\'on y custodia de evidencias}
\addcontentsline{toc}{section}{F. Procedimiento de anonimizaci\'on y custodia de evidencias}

El tratamiento de datos de los participantes durante el levantamiento de requisitos de NEXORA se rige por la Ley Org\'anica de Protecci\'on de Datos Personales (LOPDP) y su Reglamento General~\cite{lopdp2021, reglamento2023}, y se operacionaliza mediante el modelo de consentimiento informado del equipo (\texttt{PFC/E1/\allowbreak instrumentos/\allowbreak consentimiento/\allowbreak consentimiento-modelo.tex}, c\'odigo \texttt{PFC-ISR401-NEX-A03}). Para garantizar la confidencialidad y la trazabilidad verificable de las evidencias, el equipo aplica el siguiente procedimiento:

\begin{enumerate}[label=\textbf{\arabic*.}, leftmargin=1.8em, itemsep=3pt, topsep=3pt]
    \item \textbf{Asignaci\'on de c\'odigo y verificaci\'on del consentimiento:} Antes de recoger cualquier evidencia (entrevista, encuesta u observaci\'on), se asigna un c\'odigo al participante (por ejemplo, \texttt{Nexora-01}, \texttt{Nexora-02}, \texttt{Nexora-03}), se verifica que sea mayor de edad (18 a\~nos o m\'as) y se obtiene su firma voluntaria en dos ejemplares del consentimiento informado (uno para el participante y otro para el expediente restringido). La fecha del consentimiento debe ser igual o anterior a la actividad que ampara, y solo se realiza grabaci\'on de audio, video o fotograf\'ias del entorno si la fila correspondiente de la Secci\'on 7 est\'a marcada con «S\'i».
    \item \textbf{Resguardo privado del material identificable:} Los consentimientos firmados (escaneados completos, p\'aginas 1 a 3, sin cubrir datos), la carta de aceptaci\'on de la organizaci\'on cliente, las grabaciones originales y cualquier fotograf\'ia u otro archivo con datos identificables se guardan \'unicamente en la carpeta institucional privada del equipo en Google Drive (cuenta \texttt{@uteq.edu.ec}), con acceso restringido a los integrantes del equipo y al docente responsable, conforme a la Secci\'on 5 del consentimiento. Ninguno de estos archivos se sube al repositorio p\'ublico.
    \item \textbf{Limpieza y anonimizaci\'on de notas:} Las notas de campo se pasan a limpio identificando a cada participante \'unicamente por su c\'odigo y su rol gen\'erico (por ejemplo, «Profesional independiente (herrero)» o «Cliente o usuario»). Antes de publicar cualquier archivo, se revisa para retirar nombres propios, c\'edulas, firmas, tel\'efonos, correos personales, direcciones exactas, nombres identificables de negocios (salvo el de la organizaci\'on cliente, que autoriz\'o su uso por escrito) y datos de terceros que permitan reidentificar a la persona o al establecimiento.
    \item \textbf{Publicaci\'on en GitHub y registro en el manifiesto:} En \texttt{PFC/E1/evidencias/} del repositorio p\'ublico en GitHub se publican solo las notas anonimizadas y el registro de contactos sin datos personales (\texttt{registro-contacto.md}). Cada archivo privado almacenado en Drive se registra en \texttt{PFC/E1/evidencias/MANIFIESTO.md} con su nombre exacto, tipo, c\'odigo, fecha, duraci\'on cuando aplique (o «---» en documentos escaneados) y huella digital SHA-256, conservando \'integros los registros y hashes existentes.
    \item \textbf{Retiro del consentimiento y eliminaci\'on del material identificable:} Conforme a las secciones 5, 6 y 9 del consentimiento informado, el participante puede retirarse en cualquier momento entregando el desprendible de revocaci\'on o escribiendo a los correos institucionales del equipo (\texttt{lpacap@uteq.edu.ec}, \texttt{jvelezv10@uteq.edu.ec}). Ante una solicitud de retiro, se eliminan del Drive privado todos sus archivos identificables (formulario firmado, grabaciones y fotograf\'ias) y, si lo solicita y a\'un no se integraron de forma irreversible en resultados agregados, tambi\'en sus notas anonimizadas del repositorio. Con o sin retiro previo, todo el material identificable se conserva durante un plazo \textbf{m\'aximo de 24 meses desde el cierre del per\'iodo acad\'emico} y posteriormente se destruye, conservando \'unicamente los instrumentos en blanco y los resultados anonimizados.
\end{enumerate}

\textbf{Decisiones de custodia acordadas por el equipo:}
\begin{itemize}[leftmargin=1.5em, itemsep=2pt, topsep=2pt]
    \item \textbf{Eliminaci\'on de copias temporales en dispositivos personales:} Las grabaciones y escaneos almacenados en tel\'efonos o computadoras personales se eliminan apenas el archivo se sube a la carpeta privada en Google Drive y se verifica su huella SHA-256, en un plazo m\'aximo de 48 horas despu\'es de la entrevista.
    \item \textbf{Custodia y destrucci\'on de originales en papel:} Cada integrante resguarda en f\'isico los documentos firmados que recogi\'o: Paca Pilataza L. A. custodia los consentimientos de \texttt{Nexora-02} y \texttt{Nexora-03}, y V\'elez Vinces J. G. custodia el consentimiento de \texttt{Nexora-01} y la carta de aceptaci\'on. Los originales se destruyen al cumplirse los 24 meses o si un participante retira su consentimiento, dejando constancia en \texttt{MANIFIESTO.md}.
    \item \textbf{Revisi\'on cruzada previa a la publicaci\'on:} Antes de subir notas anonimizadas al repositorio, el otro integrante las revisa para confirmar que no contengan nombres, c\'edulas, tel\'efonos ni direcciones exactas; la revisi\'on se confirma por WhatsApp antes de realizar el \textit{commit}.
\end{itemize}

% ==================== G. GESTION DEL PROYECTO ====================
\section*{G. Gesti\'on del proyecto}
\addcontentsline{toc}{section}{G. Gesti\'on del proyecto}

El equipo est\'a formado por dos integrantes, por lo que cada uno asume varios roles. En esta secci\'on se presentan los roles (Tabla~\ref{tab:roles}), el cronograma hasta la semana 17 seg\'un el calendario de la asignatura (Tabla~\ref{tab:cronograma}), el plan de comunicaci\'on (Tabla~\ref{tab:comunicacion}) y los riesgos identificados con su mitigaci\'on (Tabla~\ref{tab:riesgos}) \cite{pmbok2021}.

\subsection*{G.1. Roles del equipo}
\addcontentsline{toc}{subsection}{G.1. Roles del equipo}
\begin{table}[H]
\centering
\caption{Roles y responsabilidades del equipo}
\label{tab:roles}
\footnotesize
\renewcommand{\arraystretch}{1.12}
\begin{tabularx}{\textwidth}{|L{4.2cm}|L{2.2cm}|Y|}
\hline
\rowcolor{tableheader}
\textbf{Rol} & \textbf{Responsable} & \textbf{Responsabilidades} \\
\hline
Coordinador del equipo & Jostin & Reparte las tareas, controla las fechas de cada entrega y revisa que todo se suba a tiempo. \\
\hline
Analista de requisitos & Jostin y Luis & Hacen las entrevistas, aplican la encuesta y analizan lo que sale. Jostin elabor\'o la ficha y las partes interesadas; Luis, las perspectivas y la observaci\'on. \\
\hline
Responsable de documentaci\'on (\LaTeX) & Jostin & Mantiene el \texttt{main.tex}, integra lo que escriben los dos y compila el PDF antes de cada entrega. \\
\hline
Responsable de evidencias y datos & Luis & Cuida el Drive privado, guarda los consentimientos, actualiza el \texttt{MANIFIESTO.md} y anonimiza a los participantes (Nexora-01, 02, 03, \ldots). \\
\hline
Responsable de modelos & Luis & Elabora el diagrama de contexto, los casos de uso y los diagramas UML. \\
\hline
Contacto con la organizaci\'on cliente & Jostin & Coordina las reuniones con Balanzas Electr\'onica CAS, cuyo propietario es su familiar directo (limitaci\'on declarada en la Tabla~\ref{tab:ficha}). \\
\hline
\end{tabularx}
\end{table}

\subsection*{G.2. Cronograma}
\addcontentsline{toc}{subsection}{G.2. Cronograma}
\begin{table}[H]
\centering
\caption{Cronograma del proyecto hasta la semana 17}
\label{tab:cronograma}
\footnotesize
\renewcommand{\arraystretch}{1.12}
\begin{tabularx}{\textwidth}{|C{1.4cm}|L{2.1cm}|Y|L{2.2cm}|}
\hline
\rowcolor{tableheader}
\textbf{Semana} & \textbf{Fechas} & \textbf{Actividad} & \textbf{Responsable} \\
\hline
S4 & 16--22 nov & Terminar y subir la E1: visi\'on, partes interesadas, contexto y plan de elicitaci\'on & Jostin y Luis \\
\hline
S5 & 23--27 nov & Aplicar la encuesta a usuarios en l\'inea y en persona & Jostin y Luis \\
\hline
S5 & 23--29 nov & Realizar m\'as entrevistas a profesionales independientes & Luis \\
\hline
S6 & 30 nov--6 dic & Observaci\'on de c\'omo se busca y se ofrece un servicio en Quevedo & Luis \\
\hline
S6 & 30 nov--6 dic & Tabular las respuestas de la encuesta y obtener los resultados & Jostin \\
\hline
\textbf{S7} & 7--13 dic & \textbf{E2:} evidencias, an\'alisis, priorizaci\'on y casos de uso & Jostin y Luis \\
\hline
S8 & 14--20 dic & Tutor\'ias de refuerzo y correcci\'on de lo que observe el docente en la E2 & Jostin y Luis \\
\hline
S9--S10 & 21 dic--3 ene & Redactar los requisitos funcionales y no funcionales del ERS & Jostin \\
\hline
S11--S12 & 4--17 ene & Elaborar los modelos UML y la matriz de trazabilidad & Luis \\
\hline
\textbf{S13} & 18--24 ene & \textbf{E3:} ERS seg\'un IEEE 29148, modelos UML, inspecci\'on y trazabilidad & Jostin y Luis \\
\hline
S14 & 25--31 ene & Validar los requisitos con la organizaci\'on cliente & Jostin \\
\hline
S15 & 1--7 feb & Armar el \emph{backlog} y las m\'etricas & Luis \\
\hline
\textbf{S16} & 8--13 feb & Revisi\'on final y \textbf{entrega de la E4} (se congela el s\'abado 13 de febrero, 23:59) & Jostin y Luis \\
\hline
\textbf{S17} & 15--21 feb & Preparar y presentar la \textbf{defensa} del PFC & Jostin y Luis \\
\hline
\end{tabularx}
\end{table}

\subsection*{G.3. Plan de comunicaci\'on}
\addcontentsline{toc}{subsection}{G.3. Plan de comunicaci\'on}
\begin{table}[H]
\centering
\caption{Plan de comunicaci\'on del proyecto}
\label{tab:comunicacion}
\footnotesize
\renewcommand{\arraystretch}{1.12}
\begin{tabularx}{\textwidth}{|L{3.1cm}|L{2.9cm}|L{3.2cm}|Y|}
\hline
\rowcolor{tableheader}
\textbf{Con qui\'en} & \textbf{Medio} & \textbf{Frecuencia} & \textbf{Prop\'osito} \\
\hline
Entre integrantes & WhatsApp & Todos los d\'ias & Avisar avances, resolver dudas r\'apidas y acordar qu\'e le toca a cada uno. \\
\hline
Entre integrantes & Google Meet & Una vez por semana; m\'as seguido en semana de entrega & Revisar juntos el documento y decidir lo que falta. \\
\hline
Docente (Ing. Guerrero) & Clase, tutor\'ias y correo \texttt{@uteq.edu.ec} & En cada clase; por correo ante dudas puntuales & Mostrar avances, presentar entregables y recibir correcciones. \\
\hline
Organizaci\'on cliente (Balanzas Electr\'onica CAS) & En persona o por llamada & Cada dos semanas y antes de cada entrega & Mostrar avances, confirmar requisitos y validar lo realizado. \\
\hline
Participantes (Nexora) & WhatsApp o en persona & Solo cuando se necesita una entrevista o aclarar una respuesta & Coordinar la entrevista, pedir el consentimiento y avisar si algo cambia. \\
\hline
Repositorio GitHub & \emph{Commits} & Cada vez que alguien modifica el documento & Registrar qui\'en hizo qu\'e y conservar las versiones del \texttt{main.tex}. \\
\hline
Drive privado & Carpeta compartida entre los dos integrantes y el docente & Cada vez que hay una evidencia nueva & Guardar audios, consentimientos y respuestas sin exponerlos. \\
\hline
\end{tabularx}
\end{table}

\subsection*{G.4. Riesgos}
\addcontentsline{toc}{subsection}{G.4. Riesgos}
\begin{table}[H]
\centering
\caption{Riesgos del proyecto y medidas de mitigaci\'on}
\label{tab:riesgos}
\footnotesize
\renewcommand{\arraystretch}{1.12}
\begin{tabularx}{\textwidth}{|L{4.1cm}|C{1.9cm}|C{1.3cm}|Y|}
\hline
\rowcolor{tableheader}
\textbf{Riesgo} & \textbf{Probabilidad} & \textbf{Impacto} & \textbf{Mitigaci\'on o respuesta} \\
\hline
Equipo de solo 2 personas (la gu\'ia pide de 4 a 5), con carga alta & Alta & Alto & Repartir las tareas por rol desde el inicio y adelantar trabajo en semanas con menos carga, como la S8. \\
\hline
Pocas respuestas en la encuesta & Media & Alto & Compartirla en grupos de WhatsApp de Quevedo y aplicarla tambi\'en en persona; meta m\'inima de 30 respuestas. \\
\hline
Sesgo porque la organizaci\'on cliente es familiar de un integrante & Media & Medio & Est\'a declarado en la ficha; Luis participa en las reuniones de validaci\'on y las respuestas del propietario se contrastan con las de los dem\'as participantes. \\
\hline
Conflictos en \texttt{main.tex} si ambos editan a la vez & Media & Medio & Cada uno trabaja en su secci\'on, se avisa por WhatsApp antes de editar y se hace \emph{pull} antes de cada \emph{commit}. \\
\hline
P\'erdida o exposici\'on de datos personales (audios y consentimientos) & Baja & Alto & Todo se guarda en el Drive privado; en el documento solo se usan c\'odigos Nexora y no se suben audios a GitHub. \\
\hline
Un participante retira su consentimiento & Baja & Medio & Se eliminan sus datos, se registra en el \texttt{MANIFIESTO.md} y se busca a otra persona del mismo grupo. \\
\hline
Cruce con trabajos de otras asignaturas en semanas cargadas & Alta & Medio & Revisar al inicio de cada semana la carga de las otras asignaturas y mover las tareas del proyecto a los d\'ias libres. \\
\hline
\end{tabularx}
\end{table}

% ==================== REFERENCIAS ====================
\bibliographystyle{ieeetr}
\bibliography{referencias}

\end{document}
